A deterministic supervisor can check outputs in a large language model (LLM)
pipeline and retry a stage with measured feedback. We evaluate whether the
feedback content repairs defects in a pre-registered fault-injection study of a
four-stage review-analysis pipeline. Six operators produced 142 defects over 28
workloads; the confirmatory core comprised 1,167 trials. Measured facts repaired
94.7\% of defects and a contentless rejection 94.0\% (+0.7 percentage points,
95\% CI $-2.9$ to +4.3; exact McNemar $p=0.754$). In an exploratory analysis,
measured facts increased unconditional exact restoration to the clean output by
14.6 points (68.1\% vs 53.5\%, 95\% CI +8.2 to +21.3). At the reasoning stage, restored outputs were
associated with less downstream report divergence than repaired rewrites;
these post-treatment subgroups do not establish causation. A zero thinking-token
budget retained high repair but some trials needed schema retries. A small local
model showed an exploratory +6.8-point repair contrast, with sensitivity to
repetition and model/runtime differences. Production retry was accepted on
28.2\% of real defects versus 98.5\% of injected defects; acceptance is not
correctness. Showing the defective artifact reduced measured collateral damage
from 38.4\% to 1.1\% in the placement cell. These results distinguish local
repair, restoration, verifier acceptance and downstream preservation.
